% ID: SEL-D0051-OA-E03
% ESTADO: cotejada. Segunda pasada visual y borrador compilado revisados.
% Fuente: Selectividad/sel_2013_matematicas_aplicadas/septiembre/reserva a/EXAMEN 3 -ANDALUCÍA.pdf
% SHA-256: 35f960392892cf4aa4b4a3e5b0a7b6150b323d584bb033f75db6a182b19206d4
% Página: 1; opción A; ejercicio original 3.

En una urna A hay 10 bolas verdes y 10 rojas, y en otra urna B hay 15 verdes y 5 rojas. Se lanza un dado, de forma que si sale múltiplo de 3 se extrae una bola de la urna A y en el resto de casos se extrae una bola de la urna B.
\begin{enumerate}[label=\alph*)]
\item \textbf{(1.5 puntos)} Calcule la probabilidad de que la bola extraída sea roja.
\item \textbf{(1 punto)} Si la bola extraída resulta ser de color verde, \textquestiondown{}cuál es la probabilidad de que proceda de la urna B?
\end{enumerate}
